% !TeX root = ../papers/nc-p3.tex
\section{点グラフ内の任意の平面三次印からの中心一次表示}
\label{sec:plane-cubic-mark-normality}
\begin{theorem}\label{thm:plane-cubic-mark-normality}
標準有限Frobenius窓で左右が点ファイバー、共通グラフ $X$ が一次元とする。
一位置の像が平面内に純な三次部分曲線 $C$ を含めば、全位置のその平面を消す一次式は正則正規列になる。
$C$ の重複成分、その上の埋込み点、外の余分な曲線や有限構造を許す。
全窓は定理\ref{thm:regular-normal-linear}の三生成三次元AS商の中心一次PBW持上げとなり、
全次数で標準Koszul左右AS正則、最小周期1である。
逆に三生成商の点対応が平面三次のグラフとなる同定理の全PBW持上げで、
四文字の各辺が点ファイバーかつ一次元ならこの枝に入る。
\end{theorem}
\begin{proof}
$C$ は純で $p_a=1$ であり、全位置で次数3であるから
定理\ref{thm:septic-point-graph}の残余三次の証明により各位置で平面三次になる。
その平面方程式を $t_i$、平面上の偏極を $L_i$ とする。
$H^0(\mathcal O_C)=k$、$\omega_C=\mathcal O_C$、全切断が埋込みを与える。
隣接する三切断系の積の像は六次元である。同じ束なら平面に二次式がないことから従う。
相異なるなら基点自由な鉛筆の列で $H^1(L_{i+1}/L_i)=0$ を使う。
この消滅は補題\ref{lem:genus-one-quartic-mark}と同じく、
多重次数零の束に切断があれば自明であることによる。実際、真の純部分曲線は平面で次数2以下の因子で
$H^1(\mathcal O)=0$ であり、同補題の部分曲線への双対性の証明をそのまま使える。
偏極が等しければ三次曲線に二次式がないため平面対角全体が入ってしまう。
六関係の隣接平面積への制限は次元3以下で、二式以下なら曲面が残るためちょうど3である。
この三式は平面積で完全交叉、既にこれに入る $C$ の両側次数が3なので完全交叉自体が $C$ である。

定理\ref{thm:plane-cubic-point-graphs}の記法で、平面上で零になる三関係を
$t_i B_i Y-t_{i+1} A_i X+c_i t_it_{i+1}$ と書く。
$A_i$ の核に射影点 $p$ があれば $B_{i-1}$ も特異となることを示す。
$p$ が三次曲線外なら制限三式の平面上での評価が可逆であり、次点も前点も平面外となる。
実際、混合式は次点に何も課さず、残りの三式は独立なので唯一の次点がある。
前点の存在は両側グラフの同定による。そこで一つ前の混合式で $B_{i-1}p=0$ となる。

$p\in C$ の場合は、曲線の滑らかさや、$X$ がそこで $C$ と一致することを要しない。
局所環 $\mathcal O_{X,p}$ で $C$ のイデアルを $J$ とし、余法方向の最小生成元空間
$T=J/\mathfrak m_pJ$ を全位置で同一視する。局所自明化した $t_i,t_{i+1}$ は $T$ を張り、
この二元の線型関係はちょうど三混合式の二列 $B_i j_{i+1}(p),-A_i j_i(p)$ の行空間である。
実際、平面積上の三式の持上げと二つの平面式は、周囲の滑らかな六次元空間で
$C$ のイデアルの長さ5の正則列である。極小生成元空間でまず持上げ三元を除き、
残る二元に混合三式の定数係数だけが作用する。これは $J/\mathfrak m_pJ$ の表示そのものである。
$A_i j_i(p)=0$ ゆえ、他方の列が非零なら $T$ は一次元で $t_i=0$、
他方も零なら $T$ は二次元である。前の辺の二元 $t_{i-1},t_i$ も同じ $T$ を張る。
前者の場合 $t_{i-1}\ne0$ なので混合式が $B_{i-1}j_i(p)=0$ を強制する。
後者なら二元は独立で、やはりこの列は零である。
以上どちらでも $B_{i-1}$ は特異となる。
左右核から $u\otimes t_i\in R_{i-1}$、$t_i\otimes v\in R_i$ なる非零純関係が生じ、
三次交叉に純テンソルが入る。しかし四次対合と点ファイバー条件から
どの非零三次交叉も両端支持が少なくとも3である。矛盾により全 $A_i$、逆向きに全 $B_i$ は可逆である。

商三関係 $S_i$ の非零元は階数1を持たない。もし $a\otimes b$ があれば $C$ の各一般点で積が零となる。
各平面の非零線型式は高々一つの直線成分だけを含み、その一般点の
$k(C_{\rm red})[u]/(u^m)$ で消せる重複は一層だけである。平面三次の局所式が
その直線の方程式の $m$ 乗と単元の積だからである（$m=1$ ではその一層全部）。
成分ごとの次数は次位置でも等しく、二式で消せる成分次数の合計は重複込みで高々2、三次全体は消えない。

商三次交叉の非零テンソルの第一支持が二次元以下なら、それを消す非零文字評価 $p$ を選ぶ。
末尾の全縮約は $S_i$ の $p$ での評価核に入るが、この核は点ファイバー条件から高々一次元である。
よって末尾支持が一次元となり、先頭縮約で $S_{i+1}$ に非零な階数1が生じる。従って両端支持は3である。
このようなテンソルは三文字の一端と三次元関係空間の間で同型を与えるので、
交叉は高々一次元である（代数閉体上で可逆行列だけの射影直線はない）。
他方、混合行列の可逆性から $t_i V_{i+1}=V_i t_{i+1}$ は四次元で、
三文字商の三次元数は $20-10=10$ 以上である。従って交叉はちょうど一元、商三次元数は10であり、
$t$ の二次から三次への左右乗法は単射である。一次の可逆な移送を関係に二回用いれば
この単射性で二次関係空間そのものを次へ運ぶ。従って中心表示への同定、
補題\ref{lem:ternary-finite-graph-test}による商の正則性、三次PBW条件からの全次数正則性は
定理\ref{thm:plane-cubic-point-graphs}と同じである。
逆向きは中心PBW表示を平面対に制限すればよい。
\end{proof}

\paragraph{指定を忘れた同型。}
定理\ref{thm:regular-normal-linear}の有限係数式でPBW補正と指定平面の同型を判定する。
平面印の総体は各位置での平面をGrassmann変数に取り、二つの連続する平面積上で
三式の完全交叉三次グラフとその共有中間像の一致を課して記述できる。
指定を忘れた比較には、第二窓で第一側の印に対応する印を許す。
厳密には二窓延長の系\ref{cor:two-window-extension}により、両関係を運ぶ三位置の座標変換を検査すれば十分で、
第一の印が移った後の平面はその変換で定まる。従って二重直線や三重直線を含む場合に印の個数の有限性を仮定しない。
